\chapter{First Contact: Registers, Memory, and the Call}
\label{ch:asm1}

This chapter begins the computer architecture track. By its end you will have a
running \riscv assembly program on both QEMU and, if you have one, the physical
board: a program that clears a console, prints text, and makes a sound. More
importantly, you will understand the three ideas on which everything else rests:
what a register is, how a value lives in memory, and what happens when you call a
function. We build these ideas on the smallest possible game, the ASCII version
of Pong, exactly as the platform's own tutorial does \citep{prg32tutorialascii}.

\section{The machine you are programming}

A \riscv processor of the kind in your board is, from the programmer's point of
view, astonishingly simple. It has thirty-two registers, each holding a 32-bit
value, and it executes instructions that, with few exceptions, do one of three
things: an arithmetic or logical operation between registers, a transfer of a
value between a register and memory, or a change in which instruction runs next.
That is very nearly the whole machine. The richness of computing comes not from a
large instruction set but from combining a small one in long sequences.

\begin{prgconcept}
A register is a tiny, extremely fast storage location inside the processor that
holds one 32-bit value. The processor can only compute on values that are in
registers. To work on a value stored in memory, you must first \emph{load} it
into a register; to keep a result, you must \emph{store} it back to memory.
\end{prgconcept}

The thirty-two registers have conventional roles, fixed not by the hardware but
by an agreement called the \emph{calling convention} or \emph{ABI}. Following the
convention is what lets your assembly code and the framework's C code call each
other. Table~\ref{tab:registers} shows the registers you will actually use; the
full convention is in Appendix~\ref{app:api}.

\begin{center}
\begin{tabularx}{\linewidth}{@{}llX@{}}
\toprule
Register & Name & Role in this book \\
\midrule
\reg{a0}--\reg{a7} & arguments & values passed into and out of calls; \reg{a0} also holds the return value \\
\reg{t0}--\reg{t6} & temporaries & scratch values; a call may destroy them \\
\reg{s0}--\reg{s11} & saved & values a call promises to preserve \\
\reg{ra} & return address & where a \code{ret} jumps back to \\
\reg{sp} & stack pointer & top of the stack; kept 16-byte aligned \\
\reg{zero} & zero & always reads as 0 \\
\bottomrule
\end{tabularx}
\captionof{table}{The \riscv registers you will use, and their agreed roles
\citep{prg32abi}.}
\label{tab:registers}
\end{center}

\begin{prgconcept}
The \textbf{calling convention} is an agreement about which registers carry
arguments (\reg{a0}--\reg{a7}), which carry results (\reg{a0}), and which survive
a function call (\reg{s0}--\reg{s11}, but not \reg{t0}--\reg{t6}). \prg follows the
standard \riscv convention, which is exactly why an assembly game can call a C
helper such as \code{prg32\_gfx\_clear} \citep{prg32tutorial}.
\end{prgconcept}

\section{The smallest program that runs}

Recall the three-function contract from Chapter~\ref{ch:platform}. The smallest
legal \prg game provides \code{init}, \code{update}, and \code{draw} that each do
nothing but return. In \riscv assembly, ``return'' is the instruction \code{ret},
and ``do nothing else'' is simply not writing anything else. We name our game
\code{hello} and give every symbol that prefix.

\lstinputlisting[
    style=asm,
    caption={%
        \fname{hello\_world\_asm.S} \protect\\
        The smallest \prg assembly game: three functions that return immediately.
        Build it, confirm it runs, then add to it.%
    }
]{examples/ch03/hello_world_asm.S}

A few details deserve explanation. The directive \code{.option norelax} tells the
assembler not to apply an optimisation that can interfere with the way \prg
cartridges are linked; every example begins with it \citep{prg32tutorialgraphic}.
The \code{.section .text} directive says that what follows is executable code, as
opposed to data. The \code{.global} directives make the three symbols visible to
the framework, which must find them by name. Without \code{.global}, your
functions would exist but the loader could not see them.

\begin{prgcheck}
Build this skeleton for QEMU and confirm it runs without crashing. You will see
no game yet: that is correct. The achievement is that the framework found your
three functions, called them, and survived. Appendix~\ref{app:env} has the exact
build commands for your operating system; Appendix~\ref{app:worked} shows this
skeleton wired into the firmware.
\end{prgcheck}

\section{Simple instructions and program structure}

A working program is built from a few very small operations. These
operations are provided by the RISC-V Instruction Set Architecture \citep{riscvisa}.
We will now introduce the foundational operations you need for writing
your first programs.

\subsection{Arithmetic and bit operations}

The first group is the arithmetic and bitwise instructions. In a real program
these are used to update counters, combine values, compute flags, and prepare
arguments for later work. They all read values from registers and write the
result back into a register. The result is stored in the first register.

\begin{center}
\begin{tabularx}{\linewidth}{@{}lXl@{}}
\toprule
Instruction & What it does & Example \\
\midrule
\code{add} & Add two registers & \code{add t0, t1, t2} \\
\code{sub} & Subtract one register from another & \code{sub t0, t1, t2} \\
\code{xor} & Bitwise exclusive-or of two registers & \code{xor t0, t1, t2} \\
\code{or} & Bitwise OR of two registers & \code{or t0, t1, t2} \\
\code{and} & Bitwise AND of two registers & \code{and t0, t1, t2} \\
\bottomrule
\end{tabularx}
\end{center}

The RISC-V ISA also provides immediate forms of many of these operations.
These instructions take one register operand and one small constant, which is
useful for adding a fixed amount or applying a fixed mask in one step.

\begin{center}
\begin{tabularx}{\linewidth}{@{}lXl@{}}
\toprule
Instruction & What it does & Example \\
\midrule
\code{addi} & Add a register and a constant & \code{addi t0, t1, 4} \\
\code{xori} & Exclusive-or a register with a constant & \code{xori t0, t1, 1} \\
\code{ori} & OR a register with a constant & \code{ori t0, t1, 0xff} \\
\code{andi} & AND a register with a constant & \code{andi t0, t1, 7} \\
\bottomrule
\end{tabularx}
\end{center}

\subsection{Load and store word}

A second group of instructions moves whole 32-bit values between memory and
registers. In a real program these are the only instructions that touch the
state stored in the program data; they let the program remember values
from one frame to the next.

\begin{center}
\begin{tabularx}{\linewidth}{@{}lXl@{}}
\toprule
Instruction & What it does & Example \\
\midrule
\code{lw} & Load a 32-bit word from memory into a register & \code{lw t0, 0(t1)} \\
\code{sw} & Store a 32-bit word from a register into memory & \code{sw t0, 0(t1)} \\
\bottomrule
\end{tabularx}
\end{center}

Each load or store combines a register holding a memory address with a small
constant offset. That makes it easy to access fields inside a data structure or
variables laid out in memory. You can use a zero offset when you don't need it.

\subsection{Pseudo-instructions}

The final group is the pseudo-instructions. They are not separate hardware
operations; instead, they are convenient shortcuts that the assembler expands 
into one or more real instructions. In everyday code they make assembly 
easier to read and write.

\begin{center}
\begin{tabularx}{\linewidth}{@{}lXl@{}}
\toprule
Pseudo & What it does & Example \\
\midrule
\code{la} & Load the address of a symbol into a register & \code{la a0, title} \\
\code{li} & Load a constant value into a register & \code{li t0, 42} \\
\code{mv} & Copy one register to another & \code{mv t0, t1} \\
\code{not} & Bitwise complement of a register & \code{not t0, t1} \\
\bottomrule
\end{tabularx}
\end{center}

In practice, \code{la} is used whenever a program needs the address of a string
or a data variable. \code{li} is used to load a number directly into a register
without first storing it in memory. \code{mv} is simply a copy, and \code{not}
computes the one’s complement of a register value.

With these instruction families, the RISC-V ISA provides the basic operations a
program needs: arithmetic and bitwise logic, memory load/store, and the
pseudo-instructions that make assembly manageable. The next section shows how
this instruction set is used when your assembly calls into the framework.

\section{Calling a C helper: the stack frame}
A game that only returns is not interesting for long. The first useful thing we
can do is print text, using the framework's console. The console is written in C,
so calling it means crossing the boundary between your assembly and the
framework's C, and that boundary has one rule a beginner must internalise.

When you execute \code{call}, the processor records, in the register \reg{ra},
the address it should return to. But the C function you are calling will itself
use \reg{ra} for its own calls, overwriting yours. If you do nothing to protect
your \reg{ra}, then when your function tries to \code{ret}, it will jump to the
wrong place and your program will misbehave or crash. The solution is to save
\reg{ra} on the \emph{stack} before the call and restore it afterwards.

\begin{prgconcept}
The \textbf{stack} is a region of memory that grows downward and is used to save
values across function calls. Before calling a C helper, an assembly function
lowers the stack pointer \reg{sp} to make room, stores \reg{ra} there, makes its
call, then loads \reg{ra} back and raises \reg{sp}. This save/restore pair is the
single most important habit in \riscv assembly.
\end{prgconcept}

Here is the pattern, which you will repeat in nearly every function:

\lstinputlisting[
    style=asm,
    caption={%
        \fname{print\_asm.S} \protect\\
        The stack-frame idiom. Save \reg{ra}, do work that may call C, restore \reg{ra}, return.%
    }
]{examples/ch03/print_asm.S}


Read this slowly, because it contains five distinct ideas. The \code{addi sp, sp,
-16} subtracts sixteen from the stack pointer, reserving sixteen bytes; we use a
multiple of sixteen because the convention requires \reg{sp} to stay 16-byte
aligned across calls \citep{prg32abi}. The \code{sw ra, 12(sp)} stores the word in
\reg{ra} at an offset of twelve bytes above the new stack top. The \code{la a0,
title} loads the \emph{address} of our string into \reg{a0}, because the
convention puts the first argument in \reg{a0} and the console function expects a
pointer. The two \code{call} instructions invoke C helpers. Finally the load and
\code{addi} undo the save and the reservation, in the reverse order, before
\code{ret}.

\begin{prgwarn}
Forgetting \code{sw ra} before a \code{call}, or \code{lw ra} after it, is the
classic first bug. The symptom is a function that ``returns to the wrong place'':
the program jumps somewhere unexpected, often crashing. If your assembly behaves
strangely around a call, check the \reg{ra} save/restore pair first
\citep{prg32tutorialascii}.
\end{prgwarn}

\begin{prgaside}
Why does the stack grow \emph{downward}, toward smaller addresses? It is a
convention almost as old as the stored-program computer, and it lets the stack
and the program's data grow toward each other from opposite ends of memory,
making the most of a fixed address space. Subtracting from \reg{sp} to allocate,
and adding to free, follows directly from that choice.
\end{prgaside}

\section{Assembly sections: the anatomy of a program}

When writing assembly, you must explicitly tell the assembler how to categorize different 
parts of your program using directives. The \code{.section} keyword is the switch that selects 
where the following code or data belongs: \code{.data} for mutable state, \code{.rodata} for constants, 
and \code{.text} for executable instructions. The \code{.data} section is for mutable state—variables 
like our counter that the program must modify during execution. Conversely, \code{.rodata} stands for 
``read-only data'' and is reserved for constants, such as our game title string. On physical hardware, 
placing constants in \code{.rodata} allows the toolchain to keep them safely locked in flash memory so 
they cannot be accidentally overwritten by a bug.

The very first directive in a PRG example is often \code{.option norelax}. That tells the assembler 
not to apply its relaxation pass, so it does not rewrite certain branch or symbol addresses later. 
For PRG32 cartridges, the exact layout of code and data matters, and \code{.option norelax} helps keep 
that layout predictable.

While section directives dictate \emph{where} data lives in memory, the \code{.global} directive 
dictates \emph{who} can see it. By default, every label you define in an assembly file is strictly 
private to that file. By marking a label as \code{.global}, you export it into the project's global 
symbol table. This step is what bridges your assembly code to the rest of the ecosystem, explicitly 
giving the game framework permission to locate, link against, and jump into your functions from the outside.

\section{Values in memory: the \code{.data} section}

A game must remember things between frames: where the ball is, what the score
is. In assembly, persistent values live in the \code{.data} section, each given a
label and an initial value. This is the assembly equivalent of a global variable.
For ASCII Pong, the platform's example declares its state exactly this way
\citep{prg32repo}:

\lstinputlisting[
    style=asm,
    caption={%
        \fname{game\_state\_data.S} \protect\\
        Game state in \code{.data}. Each label names a memory location holding one 32-bit word.%
    }
]{examples/ch03/game_state_data.S}

To use such a value you do two things, and the distinction between them is the
crux of understanding memory. First you load the \emph{address} of the label into
a register with \code{la}. Then you load the \emph{value} at that address with
\code{lw}, or store a new value with \code{sw}.

\lstinputlisting[
    style=asm,
    caption={%
        \fname{load\_modify\_store.S} \protect\\
        The load--modify--store cycle: read a value from memory, change it in a register, write it back.%
    }
]{examples/ch03/load_modify_store.S}

\begin{prgconcept}
A label such as \code{hello\_x} names a \emph{location} in memory. \code{la t0,
hello\_x} puts that location's address in \reg{t0}; \code{lw t1, 0(t0)} reads the
\emph{contents} of that location into \reg{t1}. The processor computes only on
contents, in registers, then writes results back. This load--modify--store cycle
is how every variable in every program ultimately works.
\end{prgconcept}

\begin{prgwarn}
Keep the two clearly apart in your mind: \reg{t0} holds an \emph{address}, \reg{t1}
holds a \emph{value}. Confusing them, using an address where a value belongs, or
vice versa, produces wrong numbers and is one of the most common beginner errors.
A useful habit, used throughout the platform's examples, is to comment each
register with whether it holds an address or a value.
\end{prgwarn}

\section{Putting it together: a console game that ticks}

We now have every piece needed for a minimal but genuine program: an \code{init}
that prints a title, an \code{update} that advances a counter each frame, and a
\code{draw} that prints the score to screen.
We also included a short beep sound that works with the \code{prg32\_audio\_beep} call.
\citep{prg32tutorialascii}.

\lstinputlisting[
    style=asm,
    caption={%
        \fname{counter.S} \protect\\
        A counter example that displays an incrementing score.%    
    }
]{examples/ch03/counter.S}

This short program demonstrates the core design pattern of every game on this platform. 
Notice how the work is split across distinct phases: \code{counter\_update} handles the logic 
by modifying our variable in memory, while \code{counter\_draw} focuses purely on presentation. 
This separation of concerns is fundamental; it ensures that your game mechanics remain entirely 
independent of how or when the screen is drawn.

The program also highlights how data types behave at the machine level through our choice 
of console helpers. While \code{prg32\_console\_write} requires the memory address of a null-terminated 
ASCII string, \code{prg32\_console\_hex32} takes a raw 32-bit numeric value directly in \reg{a0}. It 
reads the content of the register itself and formats that raw integer into human-readable text 
on the fly. Confusing these two---passing a raw number where an address is expected---is a classic 
recipe for a crash.

Finally, look at the entry and exit points of all three functions. Every single one opens 
with an identical prologue to lower \reg{sp} and save \reg{ra}, and closes with a matching 
epilogue to restore them. Because every function here uses a \code{call} instruction to hand 
off work to the framework, they would all overwrite their own return path without this defensive 
housekeeping.

\begin{prgtargets}
\textbf{Both targets.} On QEMU you will see the title on the serial monitor; the
beep is accepted but silent because the emulator does not model the buzzer. On the
ESP32\nobreakdash-C6 you will additionally hear the beep through the buzzer. The
\emph{source is identical}; only the build directory and configuration profile
differ, as shown in Chapter~\ref{ch:platform}.
\end{prgtargets}

\begin{prgtry}
Change the beep frequency from \code{440} to \code{880} and rebuild. On the board,
the tone rises by an octave. Then change the increment in \code{hello\_update}
from \code{1} to \code{2} and predict, before running, what the counter does.
Verify your prediction. This ``change one thing, predict, verify'' loop is the
method of the entire book \citep{prg32repo}.
\end{prgtry}

\begin{prgcheck}
You should now be able to build and run an assembly program on at least QEMU;
explain why \reg{ra} must be saved before a \code{call}; explain the difference
between \code{la} and \code{lw}; and describe what the \code{.data}, \code{.text},
and \code{.rodata} sections are for. If any of these is shaky, re-read the
relevant section before Chapter~\ref{ch:asm2}; everything ahead builds on them.
\end{prgcheck}
